\section{Introduction}
\label{sec:intro}

Money laundering is executed as a coordinated \emph{attempt}: a group of accounts moves money
through a fan-out, a cycle, a layered stack. Research on the AMLworld
benchmark~\citep{altman2023} is nevertheless reported almost entirely as per-transaction
classification~\citep{egressy2024,lin2024fraudgt,bilgi2024}. Group-level work is now
appearing --- learned sub-network discovery~\citep{chai2023,rico2026}, collaborative mining of
laundering groups~\citep{tian2025}, pattern-level metrics~\citep{tariq2026redirect} --- but
each paper defines its own unit, its own ground truth and its own split, and none states how a
group that straddles a split boundary is scored. Numbers from these papers cannot be set side
by side.

This paper asks a prior question: \emph{what can AMLworld support when the unit of evaluation
is the laundering attempt?} We answer it in two steps. First we fix a protocol --- a predicted
ring is a set of accounts, scored against the generator's laundering-attempt blocks by
one-to-one Jaccard matching under a stated scope rule --- and audit the benchmark under it.
Second we run a reference pipeline through the protocol and report where ring recovery is
gained and where it stops, including the results that go against our own method.

Three measurements carry the argument.

\emph{Unstated evaluation choices are larger than method effects.} Re-scoring the same runs
under a different, equally defensible rule for a ring's member set moves ring-F1 by
0.35--0.38; our largest method effect is $+0.068$.
% source: docs/FINDINGS.md (C9); results/tables/r14_scope_sensitivity.csv
The last eight days of the standard test split hold 40.7\,\% of its positives in 0.13\,\% of
its edges, and a global score quantile over that split places only 463 of 864 seeds in the two
days that carry ordinary traffic.
% source: docs/FINDINGS.md (C1; dated note on the window-matched budget)

\emph{On HI-Small, ring recovery is bounded by observability once the detector is adequate.}
Three graph neural detectors spanning test edge F1 0.596 to 0.715 give seeds-only ring-F1 0.518
to 0.551, and nine recovery attempts after the detector leave ring-F1 flat. 44\,\% of in-scope
rings have no connected in-window piece of three or more accounts, and 81.5\,\% of the missing
members of partially recovered rings sit in a piece with no in-window path to the recovered part.
% source: results/notes/final_numbers.md (detector swap); results/notes/recovery_attempts.md; docs/FINDINGS.md (component-level scope; B13)

\emph{Below that range, the detector still matters.} On HI-Medium, a detector trained only on
HI-Small, applied without retraining, gives seeds-only ring-F1 $0.396 \pm 0.015$ (five seeds)
against 0.090 for an earlier in-domain detector that differs in architecture and training data
(two seeds), and most of that gain is ranking, not budget allocation.
% source: results/notes/final_numbers.md (HI-Medium transfer)

\paragraph{Contributions.}
\begin{enumerate}
\item \textbf{A ring-level evaluation protocol for AMLworld} (Section~\ref{sec:data}): a ring
  metric under an explicit scope rule, one-to-one Jaccard matching with a deterministic
  tie-break, a ring bootstrap that keeps multiplicities, ring-folded cross-fitting for learned
  stages, and a precision-aware paired test, released as a scorer.
\item \textbf{The HI-Small bound} (Sections~\ref{sec:findings}--\ref{sec:results-ceiling}): ring
  recovery is bounded by observability, not by detector strength or reassembly, once edge F1 is
  at or above about 0.6 --- three detectors and nine recovery attempts are flat, the loss buckets
  show where rings go, and two detectors' touched rings bound what a reallocated budget can add.
\item \textbf{Transfer} (Section~\ref{sec:results-hm}): the HI-Medium result above, with the
  decomposition into ranking and allocation.
\item \textbf{An external baseline} (Section~\ref{sec:results-baselines}): TAML, an
  account-flagging method, converted to rings and scored with its protocol differences stated.
\item \textbf{A list of evaluation defects} we found in our own pipeline and in common practice,
  with their measured effects (Sections~\ref{sec:findings} and~\ref{sec:limitations}).
\end{enumerate}

We do not claim a new detection method, that transfer beats in-domain training, or that the
expansion policy helps on HI-Medium, where it was not evaluated under transfer.
